% Options for packages loaded elsewhere
\PassOptionsToPackage{unicode}{hyperref}
\PassOptionsToPackage{hyphens}{url}
%
\documentclass[
]{article}
\usepackage{amsmath,amssymb}
\usepackage{iftex}
\ifPDFTeX
  \usepackage[T1]{fontenc}
  \usepackage[utf8]{inputenc}
  \usepackage{textcomp} % provide euro and other symbols
\else % if luatex or xetex
  \usepackage{unicode-math} % this also loads fontspec
  \defaultfontfeatures{Scale=MatchLowercase}
  \defaultfontfeatures[\rmfamily]{Ligatures=TeX,Scale=1}
\fi
\usepackage{lmodern}
\ifPDFTeX\else
  % xetex/luatex font selection
\fi
% Use upquote if available, for straight quotes in verbatim environments
\IfFileExists{upquote.sty}{\usepackage{upquote}}{}
\IfFileExists{microtype.sty}{% use microtype if available
  \usepackage[]{microtype}
  \UseMicrotypeSet[protrusion]{basicmath} % disable protrusion for tt fonts
}{}
\makeatletter
\@ifundefined{KOMAClassName}{% if non-KOMA class
  \IfFileExists{parskip.sty}{%
    \usepackage{parskip}
  }{% else
    \setlength{\parindent}{0pt}
    \setlength{\parskip}{6pt plus 2pt minus 1pt}}
}{% if KOMA class
  \KOMAoptions{parskip=half}}
\makeatother
\usepackage{xcolor}
\usepackage[margin=1in]{geometry}
\usepackage{graphicx}
\makeatletter
\def\maxwidth{\ifdim\Gin@nat@width>\linewidth\linewidth\else\Gin@nat@width\fi}
\def\maxheight{\ifdim\Gin@nat@height>\textheight\textheight\else\Gin@nat@height\fi}
\makeatother
% Scale images if necessary, so that they will not overflow the page
% margins by default, and it is still possible to overwrite the defaults
% using explicit options in \includegraphics[width, height, ...]{}
\setkeys{Gin}{width=\maxwidth,height=\maxheight,keepaspectratio}
% Set default figure placement to htbp
\makeatletter
\def\fps@figure{htbp}
\makeatother
\setlength{\emergencystretch}{3em} % prevent overfull lines
\providecommand{\tightlist}{%
  \setlength{\itemsep}{0pt}\setlength{\parskip}{0pt}}
\setcounter{secnumdepth}{-\maxdimen} % remove section numbering
\usepackage{graphicx}
\usepackage{float}
\usepackage{amsmath}
\ifLuaTeX
  \usepackage{selnolig}  % disable illegal ligatures
\fi
\usepackage{bookmark}
\IfFileExists{xurl.sty}{\usepackage{xurl}}{} % add URL line breaks if available
\urlstyle{same}
\hypersetup{
  hidelinks,
  pdfcreator={LaTeX via pandoc}}

\author{}
\date{\vspace{-2.5em}}

\begin{document}

Test passed Test passed

\section{Question}\label{question}

María lleva un registro detallado de los gastos relacionados con su
vehículo. La tabla muestra los gastos semanales durante un mes completo,
organizados por categorías.

\includegraphics[width=0.9\textwidth,height=\textheight]{tabla_gastos.png}

Para analizar qué semana representó el mayor porcentaje del gasto total
mensual, resuelva paso a paso:

\textbf{IMPORTANTE - Formato de números:}

\begin{itemize}
\tightlist
\item
  \textbf{Valores monetarios}: Sin separador de miles, use punto para
  decimales

  \begin{itemize}
  \tightlist
  \item
    Ejemplo: \$16850.32 (no \$16.850,32 ni \$16850,32)
  \end{itemize}
\item
  \textbf{Respuestas numéricas}: Sin separador de miles, use punto para
  decimales

  \begin{itemize}
  \tightlist
  \item
    Ejemplo: 1234.5678 (no 1.234,5678 ni 1234,5678)
  \end{itemize}
\end{itemize}

\subsubsection{Paso 1: Identificación del mayor gasto
semanal}\label{paso-1-identificaciuxf3n-del-mayor-gasto-semanal}

Observe la tabla y identifique cuál fue el mayor gasto semanal (total
por semana).

\textbf{Respuesta:} \$\#\#ANSWER1\#\#

\subsubsection{Paso 2: Identificación del gasto total
mensual}\label{paso-2-identificaciuxf3n-del-gasto-total-mensual}

Según la tabla, ¿cuál es el gasto total del mes (suma de todas las
semanas)?

\textbf{Respuesta:} \$\#\#ANSWER2\#\#

\subsubsection{Paso 3: Configuración de la fórmula de
porcentaje}\label{paso-3-configuraciuxf3n-de-la-fuxf3rmula-de-porcentaje}

Para calcular el porcentaje, complete la fórmula con los valores
identificados:

Porcentaje = ( \#\#ANSWER3\#\# ÷ \#\#ANSWER4\#\# ) × 100\%

\subsubsection{Paso 4: Verificación de
valores}\label{paso-4-verificaciuxf3n-de-valores}

Confirme que los valores del numerador y denominador son correctos: -
Numerador (mayor gasto semanal): \$\#\#ANSWER3\#\# - Denominador (gasto
total mensual): \$\#\#ANSWER4\#\#

\subsubsection{Paso 5: Cálculo del porcentaje
final}\label{paso-5-cuxe1lculo-del-porcentaje-final}

Complete el cálculo del porcentaje:

Porcentaje = ( \#\#ANSWER3\#\# ÷ \#\#ANSWER4\#\# ) × 100\% =
\#\#ANSWER5\#\#\%

\subsubsection{Paso 6: Confirmación del tipo de gráfica (CON
PUNTUACIÓN)}\label{paso-6-confirmaciuxf3n-del-tipo-de-gruxe1fica-con-puntuaciuxf3n}

Basándose en su análisis anterior, \textbf{seleccione qué tipo de
gráfica permite identificar más fácilmente qué semana representó el
mayor porcentaje del gasto total mensual}:

\#\#ANSWER6\#\#

\textbf{Conclusión:} La semana con mayor gasto representó el
\#\#ANSWER5\#\#\% del gasto total mensual.

\subsection{Answerlist}\label{answerlist}

\begin{itemize}
\tightlist
\item
  Gráfica circular por semana
\item
  Gráfica circular por categoría
\item
  Gráfica de barras apiladas por semana
\item
  Gráfica de barras agrupadas por categoría
\end{itemize}

\section{Solution}\label{solution}

\textbf{NOTA IMPORTANTE - Configuración de evaluación automática:}

\begin{itemize}
\tightlist
\item
  \textbf{Tolerancias configuradas}: Tolerancia 1 para respuestas
  numéricas monetarias, tolerancia 0.1 para porcentajes, tolerancia 0
  para respuestas schoice
\item
  \textbf{Justificación}: Los valores monetarios son enteros grandes,
  tolerancia 1 evita rechazos incorrectos por diferencias mínimas de
  formato manteniendo precisión matemática
\item
  \textbf{Formato numérico}: Sin separador de miles, punto como
  separador decimal
\end{itemize}

\subsubsection{Análisis paso a paso del problema de gastos de
vehículo}\label{anuxe1lisis-paso-a-paso-del-problema-de-gastos-de-vehuxedculo}

Este problema de \textbf{interpretación de tablas} y \textbf{cálculo de
porcentajes} requiere un análisis secuencial que demuestre el proceso de
razonamiento matemático aplicado a contextos de gastos personales:

\textbf{NOTA IMPORTANTE - Formato de números estandarizado:}

\begin{itemize}
\tightlist
\item
  \textbf{Valores monetarios}: Sin separador de miles, use punto para
  decimales

  \begin{itemize}
  \tightlist
  \item
    Ejemplo: \$16850.32 (no \$16.850,32 ni \$16850,32)
  \end{itemize}
\item
  \textbf{Respuestas numéricas}: Sin separador de miles, punto como
  separador decimal

  \begin{itemize}
  \tightlist
  \item
    Ejemplo: 1234.5678 (no 1.234,5678 ni 1234,5678)
  \end{itemize}
\item
  \textbf{Consistencia}: Mismo formato en enunciado, opciones y
  respuestas
\end{itemize}

\subsubsection{Paso 1: Identificación correcta del mayor gasto semanal
✓}\label{paso-1-identificaciuxf3n-correcta-del-mayor-gasto-semanal}

\textbf{Respuesta correcta:} \$69448

\textbf{Análisis de gastos semanales:} - Semana 1: \$69448 - Semana 2:
\$57765 - Semana 3: \$61336 - Semana 4: \$68468

La semana 1 tuvo el mayor gasto con \$69448.

\subsubsection{Paso 2: Identificación del gasto total mensual
✓}\label{paso-2-identificaciuxf3n-del-gasto-total-mensual-1}

\textbf{Respuesta correcta:} \$257017

El total mensual se calcula sumando todos los gastos semanales:
\[\text{Total mensual} = 69448 + 57765 + 61336 + 68468 = 257017\]

\subsubsection{Paso 3: Configuración correcta de la fórmula
✓}\label{paso-3-configuraciuxf3n-correcta-de-la-fuxf3rmula}

\textbf{Respuestas correctas:} Numerador = 69448, Denominador = 257017

La fórmula de porcentaje requiere:

\begin{itemize}
\tightlist
\item
  \textbf{Numerador:} El valor observado (69448 pesos)
\item
  \textbf{Denominador:} El valor total (257017 pesos)
\end{itemize}

\subsubsection{Paso 4: Verificación de valores
✓}\label{paso-4-verificaciuxf3n-de-valores-1}

Los valores son coherentes con los datos de la tabla y representan
correctamente: - El mayor gasto semanal individual - El gasto total
mensual (suma de todas las semanas)

\subsubsection{Paso 5: Cálculo del porcentaje final
✓}\label{paso-5-cuxe1lculo-del-porcentaje-final-1}

\textbf{Respuesta correcta:} 27.0\%

\[\text{Porcentaje} = \frac{69448}{257017} \times 100\% = 27.0\%\]

\subsubsection{Paso 6: Confirmación del tipo de gráfica ✓ (CON
PUNTUACIÓN)}\label{paso-6-confirmaciuxf3n-del-tipo-de-gruxe1fica-con-puntuaciuxf3n-1}

\textbf{Opciones presentadas:}

\begin{itemize}
\tightlist
\item
  \textbf{A}: Gráfica circular por semana ← \textbf{RESPUESTA CORRECTA}
\item
  \textbf{B}: Gráfica circular por categoría
\item
  \textbf{C}: Gráfica de barras apiladas por semana
\item
  \textbf{D}: Gráfica de barras agrupadas por categoría
\end{itemize}

\textbf{Análisis de la respuesta correcta:}

``Gráfica circular por semana''

\begin{itemize}
\tightlist
\item
  Esta opción permite visualizar directamente qué semana representa el
  mayor porcentaje del total mensual
\item
  Muestra claramente las proporciones de cada semana respecto al total
\item
  Facilita la identificación inmediata de la semana con mayor
  participación porcentual
\item
  Es el tipo de gráfica más adecuado para comparar partes de un todo
\end{itemize}

\textbf{Análisis de distractores:}

\begin{itemize}
\tightlist
\item
  \textbf{Gráfica circular por categoría}: Muestra proporciones de
  categorías, no de semanas
\item
  \textbf{Gráfica de barras apiladas por semana}: Muestra composición
  pero no facilita comparación de totales
\item
  \textbf{Gráfica de barras agrupadas por categoría}: Agrupa por
  categoría, no por semana
\end{itemize}

\subsubsection{Verificación del proceso de razonamiento
completo}\label{verificaciuxf3n-del-proceso-de-razonamiento-completo}

\textbf{Datos del problema:} - Mayor gasto semanal: \$69448 (Semana 1) -
Gasto total mensual: \$257017 - Porcentaje representado: 27.0\%

\textbf{El formato híbrido con puntuación dual (cloze + schoice)
garantiza que los estudiantes:}

\textbf{Parte Analítica (Pasos 1-5):} - \textbf{Lean cuidadosamente} la
tabla para extraer datos precisos - \textbf{Identifiquen correctamente}
el mayor valor semanal y el total mensual - \textbf{Configuren
correctamente} la fórmula de porcentaje paso a paso - \textbf{Realicen
cálculos} matemáticos sin saltar etapas del proceso

\textbf{Parte de Confirmación (Paso 6):} - \textbf{Demuestren
coherencia} entre su análisis numérico y la comprensión conceptual -
\textbf{Identifiquen el tipo de gráfica} más apropiado para el análisis
requerido - \textbf{Consoliden su aprendizaje} mediante validación de
resultados

\subsubsection{Conclusión}\label{conclusiuxf3n}

La semana 1 representó el \textbf{27.0\%} del gasto total mensual de
María en su vehículo.

Esta respuesta es coherente porque: - Se basa en una lectura correcta de
la tabla - Aplica correctamente la fórmula de porcentaje - El resultado
está dentro del rango esperado (0\% a 100\%) - La gráfica seleccionada
es la más apropiada para este tipo de análisis

\textbf{Verificación adicional}: El total de gastos por categorías
(Combustible: \$117986, Estacionamiento: \$83466, Peajes: \$55565) suma
exactamente \$257017, confirmando la coherencia de los datos.

\section{Meta-information}\label{meta-information}

exname: Gastos Carro Gráficas Comparación - Análisis Secuencial Cloze
extype: cloze exsolution:
69448\textbar257017\textbar69448\textbar257017\textbar27\textbar1000
exclozetype:
num\textbar num\textbar num\textbar num\textbar num\textbar schoice
extol: 1\textbar1\textbar1\textbar1\textbar0.1\textbar0 exsection:
Estadística\textbar Interpretación de
tablas\textbar Porcentajes\textbar Análisis de datos exextra{[}Type{]}:
Cálculo exextra{[}Program{]}: R exextra{[}Language{]}: es
exextra{[}Level{]}: 2 exextra{[}Competencia{]}: Interpretación y
representación exextra{[}Componente{]}: Aleatorio y sistemas de datos
exextra{[}Contexto{]}: Laboral exextra{[}Dificultad{]}: Media

\end{document}
